% p319_1; source PDF page 319; manually verified editable table.
\begin{tikzpicture}[x=1cm,y=-1cm,inner sep=0pt,outer sep=0pt]
\path[use as bounding box] (0,0) rectangle (16.000,20.690);
\draw[line width=.4pt] (0.000,0.000) rectangle (3.850,1.515);
\node[align=center,text width=3.590cm,font={\linespread{1}\fontsize{10}{13.50}\selectfont\bfseries}] at (1.925,0.757) {以危险方法危害公共安全罪的构成要件};
\draw[line width=.4pt] (3.850,0.000) rectangle (9.900,1.515);
\node[align=center,text width=5.790cm,font={\linespread{1}\fontsize{10}{13.50}\selectfont\bfseries}] at (6.875,0.757) {要点};
\draw[line width=.4pt] (9.900,0.000) rectangle (16.000,1.515);
\node[align=center,text width=5.840cm,font={\linespread{1}\fontsize{10}{13.50}\selectfont\bfseries}] at (12.950,0.757) {范例};
\draw[line width=.4pt] (0.000,1.515) rectangle (3.850,5.030);
\node[align=center,text width=3.590cm,font={\linespread{1}\fontsize{10}{13.50}\selectfont}] at (1.925,3.272) {危险的性质};
\draw[line width=.4pt] (3.850,1.515) rectangle (9.900,5.030);
\node[align=left,text width=5.790cm,font={\linespread{1}\fontsize{10}{13.50}\selectfont}] at (6.875,3.272) {必须是具体危险，不能是抽象危险};
\draw[line width=.4pt] (9.900,1.515) rectangle (16.000,5.030);
\node[align=left,text width=5.840cm,font={\linespread{1}\fontsize{10}{13.50}\selectfont}] at (12.950,3.272) {甲盗窃公共草坪上的消防栓铜芯。法院判决甲构成以危险方法危害公共安全罪。该判决有误。盗窃行为对公共安全只产生抽象危险。甲只构成盗窃罪（2023年试题）};
\draw[line width=.4pt] (0.000,5.030) rectangle (3.850,9.645);
\node[align=center,text width=3.590cm,font={\linespread{1}\fontsize{10}{13.50}\selectfont}] at (1.925,7.338) {危险的内容};
\draw[line width=.4pt] (3.850,5.030) rectangle (9.900,9.645);
\node[align=left,text width=5.790cm,font={\linespread{1}\fontsize{10}{13.50}\selectfont}] at (6.875,7.338) {限于人身伤亡和财产损失，不包括精神损害，也即限于物质性损害，不包括非物质性损害};
\draw[line width=.4pt] (9.900,5.030) rectangle (16.000,9.645);
\node[align=left,text width=5.840cm,font={\linespread{1}\fontsize{10}{13.50}\selectfont}] at (12.950,7.338) {甲将面粉装入信封内，写上“炭疽热病菌粉”，寄给乙单位，引起乙单位恐慌。法院判决甲构成以危险方法危害公共安全罪。该判决有误。甲没有危害公共安全。甲应构成投放虚假危险物质罪，该罪的保护法益不是公共安全，而是社会管理秩序};
\draw[line width=.4pt] (0.000,9.645) rectangle (3.850,12.560);
\node[align=center,text width=3.590cm,font={\linespread{1}\fontsize{10}{13.50}\selectfont}] at (1.925,11.102) {危险的样态};
\draw[line width=.4pt] (3.850,9.645) rectangle (9.900,12.560);
\node[align=left,text width=5.790cm,font={\linespread{1}\fontsize{10}{13.50}\selectfont}] at (6.875,11.102) {公共危险是指危险\CJKunderwave{具有危及多数人的}\CJKunderwave{可能性}};
\draw[line width=.4pt] (9.900,9.645) rectangle (16.000,12.560);
\node[align=left,text width=5.840cm,font={\linespread{1}\fontsize{10}{13.50}\selectfont}] at (12.950,11.102) {甲用刀在小学门口刺伤许多学生。法院判决甲构成以危险方法危害公共安全罪。该判决有误。甲应构成数个故意伤害罪};
\draw[line width=.4pt] (0.000,12.560) rectangle (3.850,17.175);
\node[align=center,text width=3.590cm,font={\linespread{1}\fontsize{10}{13.50}\selectfont}] at (1.925,14.867) {危险的程度};
\draw[line width=.4pt] (3.850,12.560) rectangle (9.900,17.175);
\node[align=left,text width=5.790cm,font={\linespread{1}\fontsize{10}{13.50}\selectfont}] at (6.875,14.867) {危险的程度要达到与放火罪、爆炸罪、投放危险物质罪的相同程度。\\{}本罪被称为兜底罪名。错误思维是，如果行为达不到放火罪的程度，那就降一格，适用兜底罪名。实际上，本罪与放火罪是并列关系，而非高低关系};
\draw[line width=.4pt] (9.900,12.560) rectangle (16.000,17.175);
\node[align=left,text width=5.840cm,font={\linespread{1}\fontsize{10}{13.50}\selectfont}] at (12.950,14.867) {甲为了报复社会，在高速公路上点燃一个很小的火堆。法院判决，危险程度不高，不定放火罪，但定以危险方法危害公共安全罪。该判决有误。危险既然没达到放火罪的程度，也就没达到本罪程度，不构成本罪};
\draw[line width=.4pt] (0.000,17.175) rectangle (3.850,20.690);
\node[align=center,text width=3.590cm,font={\linespread{1}\fontsize{10}{13.50}\selectfont}] at (1.925,18.932) {适用顺序};
\draw[line width=.4pt] (3.850,17.175) rectangle (9.900,20.690);
\node[align=left,text width=5.790cm,font={\linespread{1}\fontsize{10}{13.50}\selectfont}] at (6.875,18.932) {本罪实际上是补充罪名。如果行为符合其他犯罪的构成要件，应优先定其他犯罪};
\draw[line width=.4pt] (9.900,17.175) rectangle (16.000,20.690);
\node[align=left,text width=5.840cm,font={\linespread{1}\fontsize{10}{13.50}\selectfont}] at (12.950,18.932) {例1，劫持火车、电车的行为，应定破坏交通工具罪。\\{}例2，盗窃机动车道路上的井盖，危害公共安全，应定破坏交通设施罪（2020年、2025年试题）};
\end{tikzpicture}
